\subsection{Ermittlung der Betriebsgrenzen einer Zenerdiode}
% Alt: Aufgabe 7
\Aufgabe{
  \label{Aufg7}
  Gegeben ist die folgende Schaltung mit einer Zenerdiode und zwei Widerständen.
  Die eingesetzte Zenerdiode besitzt eine Zenerspannung von \( U_\mathrm{Z} = \mathrm{8\,V} \).
  Die maximale Verlustleistung der Diode beträgt \( P_\mathrm{tot}=\mathrm{8\,W} \).
  Die Diode arbeitet im stabilisierenden Bereich, wenn der Strom durch sie mindestens \( \mathrm{5\,mA} \) beträgt.
  Zusätzlich sind die folgenden Werte gegeben:
  $U_\mathrm{E}=20\,\mathrm{V}\pm 10\%$ und $R_\mathrm{V} = 10\,\Omega$.\\
  Welchen minimalen und maximalen Widerstandswert darf der Lastwiderstand \( R_\mathrm{L} \) nicht unter- bzw. überschreiten, damit die Diode einerseits im stabilisierenden Bereich arbeitet und andererseits nicht überlastet wird?

  \begin{figure}[H]
    \centering
    \input{Tikz/src/FigZDiodeGrenzen.tex}\alttext{Zu sehen ist eine elektrische Schaltung mit vier Bauteilen zur Bestimmung Ermittlung der Betriebsgrenzen einer Zenerdiode. Auf der linken Seite befindet sich die Spannungsquelle \(U_\mathrm{E}\), deren Spannungspfeil nach unten zeigt. In Reihe zur Spannungsquelle sind ein Widerstand \(R\) und eine Zenerdiode \(D_\mathrm{z}\) geschaltet.  Parallel zu der Zenerdiode befindet sich ein Lastwiderstand gekennzeichnet mit \(R_\mathrm{L}\). Ein Spannungspfeil in Stromrichtung kennzeichnet die Spannung der Parralelschaltung mit Der Beschriftung \(U_\mathrm{A}\).}
    \label{fig:FigZDiodeGrenzen}
  \end{figure}
}


\Loesung{
  \begin{figure}[H]
    \centering
    \input{Tikz/src/LsgZDiodeGrenzen.tex}
    \label{fig:LsgZDiodeGrenzen}
  \end{figure}

  Gegeben sind:
  \begin{equation*}
    U_\mathrm{max} = \mathrm{22\,V} \text{ und } U_\mathrm{min} = \mathrm{18\,V}
  \end{equation*}

  \textbf{Fall 1:} Bei \( U_\mathrm{max} = \mathrm{22\,V} \)
  \begin{gather*}
    U_\mathrm{RV} = \mathrm{22\,V} - \mathrm{8\,V} = \mathrm{14\,V} \\
    I = \frac{\mathrm{14\,V}}{\mathrm{10\,\Omega}} = \mathrm{1{,}4\,A} \\
    I = I_\mathrm{D} + I_\mathrm{L} \\
    I_\mathrm{D,\ max} = \frac{\mathrm{8\,W}}{\mathrm{8\,V}} = \mathrm{1\,A} \\
    I_\mathrm{L} = \mathrm{1{,}4\,A} - \mathrm{1\,A} = \mathrm{0{,}4\,A} \\
    R_\mathrm{L} = \frac{\mathrm{8\,V}}{\mathrm{0{,}4\,A}} = \mathrm{20\,\Omega} \Rightarrow R_\mathrm{L,\ max} = \mathrm{20\,\Omega}
  \end{gather*}

  \textbf{Fall 2:} Bei \( U_\mathrm{min} = \mathrm{18\,V} \)
  \begin{gather*}
    U_\mathrm{RV} = \mathrm{18\,V} - \mathrm{8\,V} = \mathrm{10\,V} \\
    I = \frac{\mathrm{10\,V}}{\mathrm{10\,\Omega}} = \mathrm{1\,A} \\
    I_\mathrm{D,\ min} = \mathrm{5\,mA} \\
    I_\mathrm{L} = \mathrm{1\,A} - \mathrm{5\,mA} = \mathrm{0{,}995\,A} \\
    R_\mathrm{L} = \frac{\mathrm{8\,V}}{\mathrm{0{,}995\,A}} \approx \mathrm{8\,\Omega} \Rightarrow R_\mathrm{L,\ min} \approx \mathrm{8\,\Omega}
  \end{gather*}

  Siehe: Abschnitt \ref{sec:SpezielleDioden}
}